\chapter{Qualitätssicherung}

Mit der Umsetzung der Kernfunktionalität endet die reine Entwicklung nicht automatisch. Für ein didaktisches Werkzeug mit fachlicher Relevanz ist die Qualität der Simulation, der Interaktion und der Wiederholbarkeit ebenso wichtig wie die technische Funktionsfähigkeit selbst. Dieses Kapitel beschreibt deshalb die Maßnahmen zur Sicherung von Korrektheit, Bedienbarkeit und nachvollziehbarer Verifikation der Anwendung.

\section{Qualitätsziele und Prüfstrategie}

Die Qualitätsstrategie verfolgt drei Ziele: fachliche Korrektheit der Simulationslogik, robuste Bedienbarkeit der Oberfläche und reproduzierbare Ergebnisse bei identischer Eingabe. Diese Ziele sind besonders wichtig, weil die Anwendung nicht nur ein Produkt demonstriert, sondern auch als Lern- und Analyseinstrument eingesetzt werden soll. Ein Fehler in der Selektionslogik oder ein inkonsistenter Zustand in der Wiedergabe wäre für didaktische Vermittlung unmittelbar problematisch.

Aus diesem Grund wurde ein mehrstufiger Prüfansatz gewählt. Die Simulation wird auf fachliche Korrektheit getestet, die Oberfläche auf Interaktions- und Zugänglichkeitsaspekte überprüft und die Datenflüsse zwischen Szenario, Lauf und Kennzahlen auf konsistente Zustandsübergänge geprüft. Die jeweiligen Prüfmaßnahmen sind damit nicht bloß technische Nebenarbeiten, sondern Teil der fachlichen Verlässlichkeit der Anwendung.

\section{Automatisierte Verifikation der Simulationslogik}

Die Kernlogik der Anwendung wird mit den in der Projektstruktur vorhandenen Vitest-Tests überprüft. Dabei werden insbesondere die deterministischen Berechnungsschritte der Scheduling-Verfahren, die Behandlung von Zeitquanten und Präemptionen sowie die korrekte Berechnung von Lauf- und Wartezeiten validiert. Die Tests sind bewusst so angelegt, dass sie die einzelnen Algorithmik-Entscheidungen nicht endlos durch UI-Komponenten replizieren, sondern die fachliche Logik direkt auf modellierter Datenbasis prüfen.

Diese Prüfstrategie ist für die didaktische Verwendung kritisch, weil sie sicherstellt, dass die in der Gantt-Darstellung gezeigten Zustände auch im fachlichen Sinne korrekt sind. Wenn ein Prozess aufgrund eines Zeitquantums unterbrochen oder nach einer Prioritätsänderung weitergeführt wird, muss die Logik dieselbe Entscheidung auch in den Kennzahlen und im Ereignislog widerspiegeln. Genau diese Konsistenz ist durch die automatisierten Tests abgesichert.

\section{Accessibility und Bedienbarkeit}

Neben der fachlichen Korrektheit wurde die Bedienbarkeit der Oberfläche systematisch berücksichtigt. Die Projektstruktur enthält Tests für Barrierefreiheit und Interaktionslogik, die insbesondere die semantische Struktur, fokussierbare Komponenten und klare Bedienbarkeit prüfen. Die Anwendung wird damit nicht nur unter funktionalen Gesichtspunkten bewertet, sondern auch hinsichtlich ihrer Zugänglichkeit und Nutzerführung.

Ein Lernwerkzeug ist nur dann wirksam, wenn die zentrale Information akustisch und visuell verständlich vermittelt wird und Nutzerinnen und Nutzer die relevanten Steuerungselemente ohne Verwirrung bedienen können. Für diesen Zweck sind saubere Statusmeldungen, verständliche Fokusführung und konsistente Bedienkonventionen von besonderer Bedeutung.

\section{Reproduzierbarkeit und Vergleichbarkeit}

Ein weiteres Qualitätsmerkmal der Anwendung ist die Reproduzierbarkeit von Läufen. Bei identischer Eingabe und identischen Parametern muss derselbe Ablauf erzeugt werden. Diese Eigenschaft ist für Lernumgebungen besonders wichtig, da sie die Vergleichbarkeit zwischen verschiedenen Scheduling-Verfahren unterstützt. Dadurch können Prozesse, Wartezeiten und Kontextwechsel in derselben Ausgangsbasis miteinander verglichen werden, ohne dass zusätzliche, nicht nachvollziehbare Einflussfaktoren die Interpretation verfälschen.

Darüber hinaus ist die Vergleichsübersicht als Qualitätsmerkmal der Anwendung zu verstehen: Sie erlaubt die direkte Gegenüberstellung mehrerer Runs und zeigt an, an welchen Stellen sich die Verfahren systematisch unterscheiden. Die Bewertungslogik ergänzt diese Analyse um eine wertungsbasierte Perspektive, ohne die fachliche Korrektheit der Simulation zu überschreiben.

\section{Fazit zur Qualitätssicherung}

Die Qualitätssicherung der Anwendung ist daher nicht als nachträgliche Kontrolle, sondern als integraler Bestandteil der Umsetzung zu verstehen. Fachliche Korrektheit, wiederholbare Simulationsergebnisse und verständliche Benutzerführung bilden zusammen die Grundlage dafür, dass die Anwendung im didaktischen Kontext belastbar und zuverlässig eingesetzt werden kann.

Die Kombination aus deterministischer Simulationslogik, automatisierten Tests und Interaktionsprüfung schafft die Voraussetzungen für eine Anwendung, die gleichzeitig wissenschaftlich nachvollziehbar, technisch robust und im Lernprozess verständlich ist.

In der Architektur des Systems ist das Szenario-Management als Teil der Anwendungslogik umgesetzt und greift auf dieselben Normalisierungs- und Validierungsmechanismen zurück wie die Erstellung neuer Szenarien. Dadurch wird sichergestellt, dass importierte Daten denselben strukturellen Regeln unterliegen wie direkt im Generator erzeugte Daten. Die so realisierte Zustandsintegrität ist für die fachliche Nachvollziehbarkeit und die Vergleichbarkeit von Läufen von zentraler Bedeutung, da sie verhindert, dass visuelle Darstellung und Simulationskontext auseinanderlaufen.

\section{Besondere technische Herausforderungen}

Eine zentrale Herausforderung bestand in der Synchronisation von Simulationsgenauigkeit und interaktiver Bedienung. Einerseits musste die fachliche Korrektheit bei Präemption, Kontextwechsel und MLFQ-Regelwerk sichergestellt werden, andererseits sollte die Anwendung unmittelbar auf Eingaben reagieren.

Ein zweiter Schwerpunkt war die Sicherung von Reproduzierbarkeit. Dazu wurden deterministische Ausführungspfade und testbare Schnittstellen priorisiert. Automatisierte Tests prüfen sowohl Kernlogik als auch zentrale Integrationsfälle, wodurch Regressionen früh sichtbar werden \cite{vitest}.

Schliesslich ergab sich ein Trade-off zwischen Detailtiefe und Ressourcenverbrauch: Dichte Snapshot-Folgen verbessern Analyse und Wiedergabe, erhöhen jedoch Speicherbedarf. Diese Spannung wurde durch parametrisierbare Snapshot-Intervalle adressiert, sodass zwischen Genauigkeit und Laufzeitverhalten situationsgerecht balanciert werden kann.

\section{Teststrategie und Qualitätssicherung}

Die Qualitätssicherung konzentriert sich auf den Simulationskern und die Anwendungslogik. Die Tests werden mit Vitest ausgeführt und prüfen sowohl isolierte Simulationseigenschaften als auch zentrale Integrationspfade der Szenarioverwaltung \cite{vitest}. Ein Schwerpunkt liegt auf der Deterministik: Wird ein identisches Szenario zweimal ausgeführt, müssen Gesamtdauer, Prozessabschluss, Ereignisfolge und Ausführungssegmente übereinstimmen.

Darüber hinaus werden die fachlich relevanten Verhaltensweisen der Verfahren mit kleinen, gezielt konstruierten Prozessmengen geprüft. Dazu gehören der Quantum-Ablauf und die Präemption bei Round Robin, die Präemption bei einer neuen Ankunft im präemptiven LCFS-Verfahren sowie die Herabstufung eines Prozesses nach Ablauf seines Quantums in der MLFQ. Ergänzend prüft die Eingabevalidierung, ob ungültige oder unvollständige Szenariodaten normalisiert werden und ob die erzeugten Zeitsegmente gültige Grenzen besitzen.

Die Tests der Anwendungslogik ergänzen diese Kernprüfungen um die Verwaltung von Szenarien und Runs. Dabei wird unter anderem kontrolliert, ob aus einem Szenario und einem ausgewählten Run ein vollständiges Simulationsszenario erzeugt wird, ob Prozessdaten dabei kopiert werden und ob Änderungen an Szenario, Parametern oder aktivem Run zu einer veränderten Darstellung führen.

Zum Zeitpunkt der Erstellung dieser Arbeit umfasst die automatisierte Testsuite elf Testdateien mit insgesamt 34 erfolgreichen Tests. Die Tests sichern die technische Konsistenz der Implementierung ab, ersetzen jedoch keine fachliche Prüfung aller denkbaren Prozesskombinationen und keine empirische Untersuchung der Nutzung oder Lernwirkung. Die in Kapitel 7 dargestellten Produktbeispiele und Vergleichsstände dienen daher zusätzlich als nachvollziehbare Prüf- und Vergleichsfälle für die fachliche Aussagekraft der Anwendung.
